Robots are increasingly used in factories and warehouses to automate complex tasks, heavily relying on \gls{ai} and \gls{cv} to interact with their environment.
Grasping objects is a task which requires the robot to localize the object in its environment.
This environment can be quite challenging, with varying lighting conditions, reflections, dust on the camera lens, humans interacting with the robot, etc.
In such conditions, \gls{ai} systems can struggle to make accurate predictions, and fails to inform its surroundings about its uncertainty, which can lead to unsafe situations.

In this thesis, we explore where this uncertainty comes from, how to quantify it and how \gls{ai} models can be adapted and trained to deal with uncertainty.
Specifically to a regression task of predicting the center of an object in real-world situations.
Ultimately designing a model that can reliably predict the center of an object, while also providing a real-world interpretable uncertainty estimate so that actions can be taken accordingly.

We created a dataset with images of an object with ground truth labels for the position of the object in both real-world coordinates and pixel coordinates.
This dataset allows for multiple ways to balance complexity of the model and untracked data uncertainty.

We developed multiple uncertainty aware model architectures: a custom \gls{cnn} ensemble, a ResNet-18 ensemble, and a \gls{mc} Dropout model.
To validate these models, we used the \gls{nll} loss function, which balances the accuracy of the spatial predictions with the calibration of the uncertainty estimation.
After training and validating, we took the most promising models and evaluated them on a test set with a few different metrics.
The results showed that the custom \gls{cnn} ensemble performed the best, being able to predict confidently and accurately when appropriate, and successfully indicating when it was not.
However, the other models still have some improvements to be made, along with more tests on more diverse situations and different options to use the dataset.